\section{Introduction}

Breast cancer is the most commonly diagnosed cancer among women worldwide, with more than 2.3 million new cases reported annually \cite{WHO2025}. Early detection is essential for improving treatment options and patient outcome. X-ray mammography is currently the standard method for population-based breast cancer screening, but it also has relevant limitations. In particular, mammographic sensitivity is reduced in women with dense breast tissue, where the contrast between lesions and surrounding fibroglandular tissue can be limited \cite{VonEuler-Chelpin2019, Bakker2019}. In addition, mammography uses ionizing radiation and requires breast compression, which motivates continued research into complementary, non-invasive and radiation-free screening approaches.

\subsection{Background}

One possible direction for radiation-free breast cancer detection is to use the mechanical contrast between healthy and pathological tissue. Malignant breast tumours are generally stiffer than surrounding adipose and fibroglandular tissue, although reported stiffness values vary with measurement technique, loading condition and tissue type \cite{Samani2007, Ramiao2016, Goodbrake2022}. This mechanical contrast forms the basis for elastography-based imaging techniques, in which tissue deformation is used to infer differences in internal mechanical properties \cite{Patel2022, Patel2021, Jamshidi2024}.

This project was performed in collaboration with Early Warning Scan (EWS), which is developing a non-invasive screening concept based on optical and biomechanical measurements of breast deformation \cite{EarlyWarningScanWebsite}. The underlying principle is related to Digital Image Elasto-Tomography (DIET), where the breast is mechanically excited and the resulting surface motion is captured using optical measurements. Differences in internal stiffness may then affect the measured surface displacement field and dynamic response \cite{Peters2004, VanHouten2011, Fitzjohn2022}. Previous phantom and modelling studies have shown that finite element (FE) models can support the development of DIET-like systems by generating controlled deformation data and by testing the effect of stiff inclusions under known conditions \cite{Kashif2013, Ismail2018}.

Finite element method (FEM) modelling is widely used to study breast biomechanics because it allows geometry, material properties, boundary conditions and loading conditions to be controlled independently. Early breast FEM studies already demonstrated the usefulness of three-dimensional biomechanical models for simulating breast deformation and image-related applications \cite{Samani2001, Hsu2011}. More recent work has moved toward patient-specific and multi-component models that include separate anatomical regions such as adipose tissue, fibroglandular tissue, skin, muscle and support structures, each with distinct mechanical properties \cite{Samani2001, Hsu2011, Chen2024, Mazier2024}. Such models can be used to study breast motion, compression, gravity-induced deformation, image registration and the effect of internal tissue composition on the external deformation response.

However, breast modelling remains challenging. Breast tissue is heterogeneous, nonlinear and individual-specific, and its mechanical properties vary substantially between studies, measurement protocols and anatomical regions \cite{Ramiao2016, Goodbrake2022, Galbreath2025, Teixeira2023}. Anatomical features such as the posterior chestwall interface, fibroglandular distribution, skin layer and Cooper-ligament support can all influence the simulated deformation response. At the same time, adding anatomical detail increases model complexity, solver cost and the risk of poorly controlled comparisons. For an EWS-oriented modelling workflow, it is therefore important not only to make the model more realistic, but also to keep the development process reproducible and interpretable.

\subsection{Objectives}

The aim of this internship project was to develop a reproducible COMSOL-based FEM pipeline for a parametrised breast model and to use this pipeline to investigate how selected anatomical and mechanical modelling choices influence the simulated deformation response. The project started from the objective of improving the realism of an existing breast model, with an initial emphasis on geometry, posterior chestwall support and internal tissue representation. The focus was not to create a final patient-specific diagnostic model, but to establish a controlled model-development workflow in which model variants can be generated, solved and compared systematically.

The central research question of this project was:

\begin{quote}
How do controlled changes in chestwall geometry, glandular tissue distribution, outer breast shape, support representation, material settings and tumor inclusion affect the simulated displacement and stress response of a heterogeneous finite element breast model?
\end{quote}

To address this question, the model was developed in a staged manner. The first stages focused on establishing a stable dynamic baseline and improving the anatomical description of the model through the posterior chestwall geometry and glandular tissue distribution. Subsequent stages introduced matched controls for outer-shape asymmetry and Cooper-ligament support, followed by material and skin-layer sensitivity cases and an initial tumor screening route. This staged structure was chosen to make the effect of individual modelling assumptions easier to interpret.

COMSOL Multiphysics was used as the intended modelling environment for this project, but an important part of the work was to determine whether a sufficiently automated and reproducible COMSOL workflow could be built for this type of parametrised breast model \cite{COMSOLMultiphysicsWebsite}. The resulting pipeline supports systematic comparison between model variants while keeping the model definitions and evaluation outputs traceable. The final outcome of the project is therefore a reproducible COMSOL FEM model-development pipeline for EWS-oriented breast deformation studies, together with a critical assessment of which model components are suitable for quantitative interpretation and which should still be considered exploratory model-development results.

